\chapter{Autentikáció, session és jogosultságkezelés}

\noindent\textbf{Tanulási út: Gyors út: alap.} Válaszd külön az autentikációt, route policyt és objektumszintű authorizationt.\par\medskip

\rustbridgeblockhu{Az enumok, typed azonosítók, \code{Result}, elágazások és normál domain predikátumok továbbra is alkalmasak policylogikára.}{A framework trusted principalt, route-szintű auth policyt és explicit authorization határt ad a védett objektumokhoz/effectekhez.}{Ne cseréld le a domain rule-okat framework-szókincsre. Az authentication identitást ad; a typed business rule továbbra is authorityt dönt, a Velran pedig azt kényszeríti ki, hol kell ennek bizonyítva lennie.}


\invariantblock{Biztonsági invariáns: az identitás nem jogosultság}{Az authentication azt állapítja meg, ki a hívó. Az authorization azt, hogy ez a hívó elvégezheti-e ezt a műveletet ezen az objektumon. Az authorization a protected read vagy mutation előtt fejeződjön be.}

\diagnosticblockhu{Verifier diagnostic}{SEC-A01-022}{A route identitása authenticated, de a role/permission proof nem elégíti ki a handler követelményét.}{Erősítsd a route policyt, vagy csak akkor módosítsd a handler követelményét, ha azt az üzleti authorization contract valóban megengedi.}

Az authentication identitást bizonyít; az authorization azt dönti el, mit tehet a hívó. A session a trusted identityt viszi tovább a kérések között. A Velranban az alkalmazás route-policyja \code{.vrn}-ban, az identity- és session-policy pedig szerverkonfigurációban marad.

\section{Hol helyezkedik el az authentication és authorization az életciklusban?}

A request-életciklus segít különválasztani az identity és az objektum-authority kérdéseket:

\begin{lstlisting}
request
  -> parse / validate
  -> authenticate: who is calling?
  -> authorize: may caller act on this object?
  -> query / transaction
  -> response
\end{lstlisting}

A route policy a széles belépési kérdést válaszolja meg. Object-, tenant-, permission-, MFA- vagy critical-operation ellenőrzés ezt tovább szűkítheti: authenticated azt jelenti, hogy van trusted principal, nem azt, hogy a művelet engedélyezett.

A login authenticationt hoz létre: credential form és CSRF ellenőrzés, identity backend, session rotation, audit, redirect. A credential-policy az identity backend felelőssége, nem az objektumszintű authorizationé.

\section{A teljes login útvonal}

A beépített HTML login endpoint:

\begin{lstlisting}
GET  /__velran/auth/login
POST /__velran/auth/login
POST /__velran/auth/logout
\end{lstlisting}

Ezek szerver-endpointok, nem az alkalmazás által definiált route-ok. A GET login oldal egy sessionhöz kötött CSRF tokent ad. A POST login form pontosan a \code{_csrf}, \code{username}, \code{password} és \code{totp} mezőket várja.

A sikeres login magas szinten:

\begin{enumerate}
  \item ellenőrzi a form szerkezetét és a CSRF tokent;
  \item kanonizálja a felhasználónevet;
  \item alkalmazza a login rate limitet;
  \item ellenőrzi a local-auth vagy LDAP credentialt;
  \item ha szükséges, ellenőrzi a TOTP-t vagy local auth esetén a recovery kódot;
  \item új, autentikált sessionre rotál;
  \item audit eseményt ír;
  \item \code{303 See Other} válasszal a gyökérre irányít.
\end{enumerate}

A session rotation fontos: sikeres autentikáció után nem marad használatban az anonim session azonosítója. Ez session-fixation elleni alapvédelmi elem.

\section{Local auth: külön identity adatbázis}

Single-node alkalmazásnál használható a local auth SQLite store. Ez szándékosan nem az application DB. A Velran program typed SQL queryjei így nem férnek hozzá a credential táblákhoz.

Secret file például:

\begin{lstlisting}
/run/secrets/velran/local-auth-db-url
\end{lstlisting}

Tartalma:

\begin{lstlisting}
sqlite:///srv/velran/auth/users.db?mode=rwc
\end{lstlisting}

Inicializálás:

\begin{lstlisting}[language=bash]
/usr/local/bin/velran-cli auth init \
  --db-url-file /run/secrets/velran/local-auth-db-url
\end{lstlisting}

User létrehozásnál a jelszó nem CLI argumentumként szerepel, hanem fájlból olvasódik:

\begin{lstlisting}[language=bash]
printf '%s\n' 'a-very-long-random-password' \
  > /run/secrets/velran/new-user-password
chmod 600 /run/secrets/velran/new-user-password

/usr/local/bin/velran-cli auth user-add \
  --db-url-file /run/secrets/velran/local-auth-db-url \
  --username alice \
  --password-file /run/secrets/velran/new-user-password \
  --role User \
  --role Editor
\end{lstlisting}

A local store Argon2id jelszóhash-t használ, egyedi random salttal. A jelenlegi implementáció legalább 12 és legfeljebb 1024 karakteres jelszót fogad el. Ezt ne keverd össze alkalmazásoldali form-validációval: a credential-policy az identity backend felelőssége.

\section{Account lifecycle operátori CLI-val}

A felhasználó állapotát nem kell közvetlen SQL-lel módosítani. A CLI erre explicit műveleteket ad:

\begin{lstlisting}[language=bash]
velran-cli auth password-set --db-url-file ... \
  --username alice --password-file ...

velran-cli auth disable   --db-url-file ... --username alice
velran-cli auth enable    --db-url-file ... --username alice
velran-cli auth roles-set --db-url-file ... --username alice \
  --role User --role Publisher
\end{lstlisting}

A disable nem törli a usert. Kívülről a sikertelen belépés nem árulja el, hogy rossz jelszó, nem létező user vagy tiltott account áll-e a háttérben: az ilyen különbségek account-enumeration csatornát nyithatnának.

\section{Session invalidation: auth\_generation}

Local auth esetén credential- vagy authorization-változás növeli az \code{auth_generation} értékét. A session a létrehozáskori generációt tárolja; eltérésnél a régi session érvénytelen és új anonim session készül.

Ezért például az alábbi operátori műveletek nem csak a következő loginra hatnak:

\begin{itemize}
  \item password change;
  \item role change;
  \item account disable/enable;
  \item TOTP enrollment vagy disable.
\end{itemize}

Ez különösen fontos role-visszavonásnál: nem kell megvárni a korábbi session természetes lejáratát.

\section{Session store: memory vagy Redis}

Development/test környezetben használható in-memory session store. Productionban, különösen több server instance esetén, Redis legyen a shared session backend.

\begin{lstlisting}
[redis]
url_file = "/run/secrets/velran/redis-url"
allow_insecure = false
\end{lstlisting}

Redis a Velranban több rövid életű infrastruktúra-állapotot is hordozhat:

\begin{itemize}
  \item shared session;
  \item TOTP replay védelem;
  \item login rate limiting;
  \item public cache.
\end{itemize}

Redis nem üzleti source of truth. Disaster recovery után elfogadható az üres session/cache state: a felhasználók újra belépnek, a cache újraépül.

\section{Cookie policy}

Normál production módban a session cookie neve:

\begin{lstlisting}
__Host-velran_session
\end{lstlisting}

A szerver \code{Path=/}, \code{HttpOnly}, \code{Secure} attribútumokat használ. Alapesetben \code{SameSite=Lax}; credentialed cross-origin konfiguráció esetén \code{SameSite=None; Secure}.

Developmentben az explicit \code{--insecure-dev-cookies} escape hatch miatt a cookie neve \code{velran_session}, és nincs \code{Secure}. Ezt internetre kitett production deploymentben ne használd.

\section{Logout}

A logout szándékosan POST művelet:

\begin{lstlisting}[language=Velran]
<form method="post" action="/__velran/auth/logout">
    <input type="hidden" name="_csrf" value="{{ csrfToken }}">
    <button type="submit">Logout</button>
</form>
\end{lstlisting}

A szerver ellenőrzi a CSRF tokent, invalidálja a régi sessiont, létrehoz egy friss anonim sessiont, új cookie-t küld, majd auditálja a logout eseményt. Így a kijelentkezés nem egyszerű kliensoldali cookie-törlés.

\section{LDAP / LDAPS vállalati identityhez}

Shared vagy vállalati identity esetén használható LDAP backend. Productionban LDAPS URL-t használj, a service bind credential pedig secret file-ból jöjjön.

Példa konfiguráció:

\begin{lstlisting}
[auth]
ldap_url = "ldaps://ldap.example.com:636"
ldap_search_base = "ou=people,dc=example,dc=com"
ldap_username_attribute = "uid"
ldap_service_bind_dn_file = "/run/secrets/velran/ldap-bind-dn"
ldap_service_bind_password_file = "/run/secrets/velran/ldap-bind-password"
require_totp = true
\end{lstlisting}

Local auth és LDAP ne legyen egyszerre elsődleges identity backend ugyanabban a processben. Ezzel nincs bizonytalan backend-prioritás vagy azonos username-ek közötti névütközés.

LDAP használatakor a role mapping trusted szerveroldali konfigurációból jöjjön. A kliens által küldött \code{role}, \code{username} vagy hasonló mező soha nem authorization authority.

\section{TOTP és recovery code}

Local auth userhez TOTP enrollment:

\begin{lstlisting}[language=bash]
/usr/local/bin/velran-cli auth totp-enroll \
  --db-url-file /run/secrets/velran/local-auth-db-url \
  --username alice
\end{lstlisting}

A CLI megadja az enrollmenthez szükséges secret/URI adatot és recovery kódokat. A recovery kódok a local auth DB-ben csak hashként maradnak meg, és egy sikeresen felhasznált recovery kód elfogy.

TOTP kikapcsolása:

\begin{lstlisting}[language=bash]
/usr/local/bin/velran-cli auth totp-disable \
  --db-url-file /run/secrets/velran/local-auth-db-url \
  --username alice
\end{lstlisting}

Ha \code{require_totp = true}, az a user sem léphet be, akinek egyébként helyes a jelszava, de nincs használható második faktora.

A login form \code{totp} mezőjébe hatjegyű TOTP vagy local auth esetén recovery code kerülhet. Redis backend mellett a TOTP replay állapot shared, ezért több server instance is ugyanazt a replay-védelmet látja.

\section{Route authorization}

Az alkalmazás route szinten három gyakori policyt deklarálhat:

\begin{lstlisting}[language=Velran]
route account GET "/account" auth user => account;
route admin GET "/admin" auth role Admin => admin;
route security GET "/security" auth mfa => security;
\end{lstlisting}

Jelentésük:

\begin{itemize}
  \item \code{auth user}: bármely autentikált principal;
  \item \code{auth role Admin}: autentikált principal a trusted sessionben szereplő \code{Admin} role-lal;
  \item \code{auth mfa}: autentikált principal sikeresen igazolt második faktorral.
\end{itemize}

HTML route-nál a nem autentikált kérés a login oldalra irányítható. JSON API-nál ugyanez nem HTML redirectként jelenik meg, hanem \code{401 Unauthorized} válaszként; ez API kliensnél kiszámíthatóbb viselkedés.

\section{A principal használata az oldalon}

Az autholt oldal trusted sessionadatot kap:

\begin{lstlisting}[language=Velran]
#[page]
fn account(ctx: PageContext) -> Result<Html, PageError> {
    return Ok(html {
        <main>
            <h1>Account</h1>
            <p>Signed in as {{ authPrincipal }}</p>
            <form method="post" action="/__velran/auth/logout">
                <input type="hidden" name="_csrf" value="{{ csrfToken }}">
                <button type="submit">Logout</button>
            </form>
        </main>
    });
}

route account GET "/account" auth user => account;
\end{lstlisting}

A \code{authPrincipal}, a session role-jai és az MFA állapot szerveroldali trusted state. Ne próbáld ezeket hidden form mezőből vagy query parameterből visszaépíteni.

\section{Route role nem mindig elég}

A route-level role azt mondja meg, hogy a felhasználó beléphet-e egy műveletbe. Objektumszintű authorizationnál ennél több kellhet.

Példa: egy szerkesztő csak a saját draftját módosíthatja. Ilyenkor nem elég az, hogy \code{auth role Editor}: az actionnek a DB-ből olvasott objektum tulajdonosát is össze kell vetnie a trusted session principaljával.

Az alapelv:

\begin{center}
\textbf{identity a sessionből, objektumállapot a DB-ből, döntés szerveroldalon}
\end{center}

A kliens által küldött \code{owner}, \code{actor} vagy \code{role} mező legfeljebb adat, nem authority.

\subsection*{Az object authorization typed guard}

Az \code{authorize} guard csak már betöltött, \textbf{nem optional} model recordon használható. A compiler azt is ellenőrzi, hogy az owner mező létezik és \code{String} típusú. Ezért az \code{Option<Article>} értéket visszaadó query eredményét előbb explicit módon kezelni kell; authorization nem értelmez ``nincs rekord'' állapotot.

\begin{lstlisting}[language=Velran]
let article = loadArticle(db, id)?;
authorize article owner authorUsername or role Publisher;
\end{lstlisting}

Az ilyen guardot tartalmazó handler route-ja nem lehet public, és public cache sem kerülhet object authorization elé. A tulajdonos vagy a felsorolt trusted role engedélyez; minden más autentikált actor \code{403 Forbidden} választ kap.

\section{Login rate limiting}

A login külön rate-limit policyt kap. Például:

\begin{lstlisting}
[auth]
login_max_attempts = 8
login_principal_max_attempts = 16
login_source_max_attempts = 64
mfa_max_attempts = 5
login_window_secs = 300
\end{lstlisting}

A szerver a trusted proxy szabályok után meghatározott effective client IP-t és a kanonikus usernevet használja a limiter kulcsához. Ezért reverse proxy mögött a \code{trusted_proxy_cidrs} auth security szempontból is fontos.

Sikertelen credential vagy második faktor nem ad részletes hibát a kliensnek. Rate-limit esetén \code{429 Too Many Requests}, auth backend elérhetetlenségnél \code{503 Service Unavailable} a megfelelő kategória.

A limiter source+principal, principal több source-on és source több principalon dimenziót is követ; MFA/recovery esetén szigorúbb stage érvényes. Ez lefedi a distributed guessinget és credential stuffingot. Shared limiter-store hiba fail closed.

\section{CSRF, Origin és Fetch Metadata}

Session-cookie alapú autentikációnál a böngésző automatikusan elküldi a cookie-t, ezért state-changing kérésnél CSRF-védelem kell. A Velran defense-in-depth modellje több réteget használ:

\begin{itemize}
  \item session-bound CSRF token;
  \item same-origin Origin/Referer policy;
  \item Fetch Metadata kontroll;
  \item szigorú request parsing;
  \item trusted proxy policy.
\end{itemize}

A CORS nem helyettesíti a CSRF-et. Ha külön frontend originhez credentialed CORS-t engedélyezel, az külön security policy és HTTPS-et igényel.

\section{Auth audit}

A login/logout események audit logba kerülnek. Hasznos mezők például a request ID, actor, outcome és effective client IP. Az audit célja eseményrekonstrukció, nem credential-debugging.

Soha ne logold:

\begin{itemize}
  \item a jelszót;
  \item TOTP secretet vagy recovery code-ot;
  \item session ID-t vagy teljes Cookie headert;
  \item CSRF tokent;
  \item LDAP bind vagy DB/Redis credentialt.
\end{itemize}

A login eredménye lehet például siker, invalid credential, invalid second factor vagy rate limited; a kliensnek mutatott üzenet ettől szándékosan lehet kevésbé részletes.

\section{Local auth vagy LDAP?}

\begin{center}
\begin{tabularx}{\textwidth}{@{}lYY@{}}
\toprule
 & Local auth & LDAP/LDAPS \\
\midrule
Identity store & külön SQLite DB & vállalati directory \\
Tipikus deployment & single-node & shared/vállalati identity \\
User admin & Velran CLI & directory admin workflow \\
Role forrás & local auth store & trusted mapping \\
TOTP & local enrollment + recovery & server-side TOTP policy/secrets \\
Multi-instance identity & nem ez a cél & igen \\
\bottomrule
\end{tabularx}
\end{center}

Ha az alkalmazásnak saját publikus regisztráció, emailes password reset vagy önkiszolgáló MFA enrollment kell, azt ne próbáld a meglévő operátori CLI köré vékony route-tal odabarkácsolni. Ezek token-életciklust, email securityt, rate limitet és külön recovery threat modelt igényelnek.

\section{Rich text és autentikált CMS}

Az autholt CMS-szerkesztői felület sem indokolja raw HTML tárolását. A szerkesztői jogosultság és a content sanitization két külön védelmi réteg. Tartalomként Markdownot tárolj, megjelenítéskor pedig a biztonságos renderert használd:

\begin{lstlisting}[language=Velran]
<article class="article-body">
    {{ markdown_html }}
</article>
\end{lstlisting}

A renderer Velran forrásmodul, nem engine-direktíva. Típusos \code{SafeHtml} értéket állít elő; a raw HTML escape-elve jelenik meg, a linkcélokat pedig a modul és az általános SafeHtml boundary korlátozza.

\section{Production ellenőrzőlista}

Mielőtt autentikációt engedsz productionben, ellenőrizd:

\begin{itemize}
  \item pontosan egy elsődleges identity backend aktív;
  \item local auth DB elkülönül az application DB-től;
  \item jelszó és service credential csak secret file-ból jön;
  \item production sessionhöz Redis van, ha több server instance fut;
  \item a cookie Secure/HttpOnly és nincs \code{--insecure-dev-cookies};
  \item a \code{trusted_proxy_cidrs} csak a valódi Apache/Nginx proxy címét fedi;
  \item a login limiter be van állítva;
  \item a state-changing formok CSRF-et használnak;
  \item role helyett objektumszintű authorization is van, ahol szükséges;
  \item credential vagy session secret nem kerül logba;
  \item account disable/role/TOTP változás session-visszavonását tesztelted.
\end{itemize}

\section{Ellenőrzött companion példák}
A canonical futtatható companion példák:
\begin{itemize}
  \item \path{examples/auth/app.vrn} -- route-szintű authentikáció és role policy;
  \item \path{examples/object-authorization/app.vrn} -- objektumszintű authority;
  \item \path{examples/mfa-elevation/main.vrn} -- MFA elevation;
  \item \path{examples/critical-operations/main.vrn} -- critical-operation contractok.
\end{itemize}
Így ezek az authority-mechanizmusok külön maradnak, nem mosódnak egyetlen általános ``auth'' fogalommá.

\section{Gyakorlat: készíts autorizációs mátrixot}
\paragraph{Gyakorlási mód: támogatott.} Használd az előszóban meghatározott támogatott módszert.

Készíts kis táblázatot anonymous felhasználóval, hitelesített tulajdonossal, hitelesített nem tulajdonossal, privilegizált role-lal és letiltott accounttal. Sorold fel, melyik actor mit tehet egy jegyzettel. A táblázattal különítsd el a route-szintű access policy-t az objektumszintű autorizációtól.

\section{Tipikus hiba: az authenticationt authorizationnek tekintjük}
Az, hogy tudjuk, ki a hívó, nem bizonyítja, hogy hozzáférhet egy konkrét rekordhoz. A route policy adja a széles belépési szabályt; ownership-, tenant- vagy domain-specifikus ellenőrzést továbbra is a betöltött objektumon kell elvégezni.

\section{Folyamatos mintaprojekt: Secure Notes}
Használd az \path{examples/secure-notes/checkpoints/07-auth/main.vrn} snapshotot. Minden jegyzetművelethez követelj authenticationt. View, edit és delete útvonalon kényszerítsd ki az ownership szabályt. Ezután szándékosan teszteld a negatív esetet: egy érvényesen belépett felhasználó se férjen hozzá más jegyzetéhez pusztán egy azonosító kitalálásával.
